\section{Discussion of research results and decision to narrow down}
\label{sec:discussion-reduce-scope}

\subsection{Warum trotz bedenken an dem proaktiven ansatz festeglten wurde.}

Diese Arbeit ist von einer bestimmten Prämisse ausgegangen: ein KI-System, das in einem
zielorientierten Remote-Meeting moderierend eingreift, solange das Meeting läuft, und das dies
proaktiv tut.

Da faciliattorien eben falls proaktiv in meeting eingreifne und moderation teil ihrer rolle ist
fragt, Forschungsfrage~2 fragt deshalb nicht, ob Moderation grundsätzlich nützlich wäre. Sie fragt,
welche der in RQ1 beschriebenen Herausforderungen ein solches System im laufenden Meeting
überhaupt erreichen kann, und wie der Eingriff beschaffen sein muss.

Es könnte kritisert, werden dass diese Annahme zu eng war. In der Literatur wurden, beretis tendenzen
festgestellt, sich Meeting tielenhemden eher zu weniger proaktiv eingreifender systeme tendieren. (vgl. 2.3.2)

Auch aus den Experten interview gehen einige vermutete Limitation bei eingereifder KI hervor, welche exakte
solche proaktiven ansätze welche die arbeit eingetlich untersucht als nicht viel versprechend deklaieren.

Gleichzeitig geht aus der Literatur und  Interviews hervor, dassproaktives einschreiten in einigen
entscheidungs fokusierten Meeting mit convergierendem charakter, nötig und bei menschlichen Facilitatorn
gänige Praxis ist. Vom autor gesichter forschungsansätz verfolgen bei KI-Systemen aber eher zurückhalender Ansätze,
was z.B. unteranderem zur folge hat das die beiträge der KI ignoriert wurden. (vgl. 2.32)

Desweiteren sind die limitation welche die expertern interviews implizieren, auch wenn sie plausible
erscheinen, hypotetisch.

Offen bleibt die frage wie tielnehmenden in ziel-orietneiren-remote meetings eine proactives eingreifne eines
KI-Systems den tatäschlich wahrnehmen würden. Aber dem autor in seiner
literatru recherche keine systeme aufgefallen, welche tatäschlich versucht hätten in einer live-meeting
situation proaktiv einzugreifen, weshlab auch wenige system design wurden, die dies versuchen.

Daher hat sich der Atuor dazu entschieden an der ursprüglichen premisse verzuhalten und zu versuchen einen
ersten KI protoyp zu entferne welcher tatsälich in einer live meeting situation evaluiert werden kann.
Dies geschah nicht weil die geäußerten Bedenken oder Limitationen nicht für plausible Emfindungen wurden, sondern
viel mehr aus der annahme heraus, dass ein erster verusch ein solches proaktives system zu desingen, weiter
erkenntisse und herausfoderungen sowie chancen offenlegen würde, die sonst wenig beachtung gefunden hätten.
Aus diesen Leranings wird sich erfhoft patter abzuleiten werden, um solche proaktiv moderirne system weieter
zuentwickeln.

% ---------------------------------------------------------------------------------------------------------------------
\subsection{Warum eine weitere Einschränkung notwendig war}

Die vorangegangenen Kapitel öffnen einen weiten Problemraum. Zielorientierte Remote-Meetings
scheitern an Tempo, Fokus, ausbleibenden Entscheidungen, ungleicher Beteiligung und verzerrtem
Informationsaustausch (\autoref{sec:facilitation-and-group-decisions},
\autoref{sec:hidden-profiles}).

Die Experteninterviews identifizieren darin mehrere Faciliating verhaltenweise, die je nach Meeting situation
praktikabel ist. Ein system zu entwicklen, welches auf alle diese verhaltenweise reagieren kann, wäre zu
komplex undwürde RQ2 nicht beantworten. Eine beobachtete Wirkung ließe
sich keiner einzelnen Designentscheidung zurechnen.

Einige der vorschläge der Expterne fpr Potenziale des meetings fallen außerhalb des scopes. Vorbereitung und
Dokumentation (F8, F7) liegen vor und nach dem Meeting.
Zusammenfassung, Clustering und Zeit-Erinnerungen (F1, F2) nähern bestehende Werkzeuge bereits
an (\autoref{sec:existing-meeting-ai-tools}). Fehlt eine Rolle, fragen Facilitatoren die Gruppe,
was diese Rolle sagen würde. Sie wie in (F6) zu simulieren ein andere Ansazt, welcher nicht unter eine
proactive moderation fällt.  ist ein anderer Eingriff.
Ein Entscheidungsverfahren (F3) kann die Gruppe einschalten, sobald Optionen sichtbar sind. Diese
Züge bleiben gültige Praxis. Sie verlangen kein System, das ungefragt in das laufende Gespräch
eintritt.

Daher wurde als vielversprechender Ansatz der Bias-Hinweis (F5) gewählt, welcher am ehesten mit der
Froschugnsfrage einhergeht. Gleichzeitig bleibt es trotz der vermuteten Grenzen gerechtfertigt, diesen
proaktiveren Ansatz zu verfolgen, denn hier ist der Nutzen eines proaktiven Systems am größten. Da die Gruppe die
Verzerrung ggf. nicht selber erlebt, welche aktiv um hilfe geben werden müssen vermutlich gar nicht erst
getriggert und zu passive ansätze könnten ggf. ignoriert werden.

Der Bias, der an einer Kernaufgabe von Facilitatoren zusamenhängt, ist der Shared-Information-Bias
(\autoref{sec:hidden-profiles}). Teilnehmende zu integrieren heißt in dieser Praxis, die Stimme
oder den Fakt einzubringen, der das Ergebnis verändern würde
(\autoref{sec:interview-facilitation-practice}). Für Entscheidungen ist das besonders erfoderlich.
Im Hidden Profile liegt die relevante Information bereits vor, nur über die Gruppe
verteilt. Jede Person bildet auf ihrem Teilwissen eine schlechtere Präferenz, als die Gruppe sie
aus der gemeinsamen Lage bilden würde. Die bessere Option muss in der konvergenten Phase gesagt
werden. Sie muss weder beschafft noch neu erzeugt werden.

Das Inner-Thoughts-Framework biete eine technische grundlage welche durch adaption an diesen usecase
angepasst werden kann.

somit wurde diese fokusierung als sowohl viel verpsrehcen im kontext der bisherh gewonne erkennist es als
auch als technisch umsetzbar eingestuft.
